\documentclass[11pt,a4paper]{article}

% --- Packages ---
\usepackage[utf8]{inputenc}
\usepackage[T1]{fontenc}
\usepackage{lmodern}
\usepackage[margin=2.5cm]{geometry}
\usepackage{graphicx}
\usepackage{booktabs}
\usepackage{longtable}
\usepackage{enumitem}
\usepackage{amsmath}
\usepackage{hyperref}
\usepackage{xcolor}
\usepackage{natbib}
\usepackage{parskip}
\usepackage{titlesec}
\usepackage{abstract}

\hypersetup{
    colorlinks=true,
    linkcolor=blue!60!black,
    citecolor=blue!60!black,
    urlcolor=blue!60!black,
}

\titleformat{\section}{\large\bfseries}{\thesection}{1em}{}
\titleformat{\subsection}{\normalsize\bfseries}{\thesubsection}{1em}{}

\renewcommand{\abstractnamefont}{\normalfont\bfseries}
\renewcommand{\abstracttextfont}{\normalfont\small}

% --- Document ---
\begin{document}

\begin{center}
    {\LARGE\bfseries §§The Forge Protocol v2}\\[0.4cm]
    {\large An Anti-Deskilling Interaction Paradigm\\for AI-Augmented Knowledge Work}\\[0.8cm]
    {\normalsize Lorenzo Famiglini}\\[0.2cm]
    {\small \url{https://github.com/lorenzofamiglini/forge-protocol}}\\[0.6cm]
    {\small April 2026}
\end{center}

\begin{abstract}
A convergence of empirical research from 2024--2026 confirms that heavy LLM use
degrades critical thinking, weakens neural connectivity during writing, reduces
knowledge retention, and causes measurable clinical deskilling. The failure mode
is well-characterised: users dump fragments into an LLM, accept completions, and
outsource cognitive work---a textbook instance of the generation effect in
reverse. This paper presents the Forge Protocol, a structured interaction
framework grounded in cognitive science that defines four modes of LLM
interaction (Forge, Anvil, Crucible, Executor), each with explicit rules about
what the LLM does, what the user does, and when each mode applies. The protocol
is designed to reverse AI-induced deskilling by exploiting the same cognitive
mechanisms---generation effects, desirable difficulties, cognitive forcing
strategies---that the standard LLM interaction pattern destroys. Design
principles draw extensively on Cabitza et al.'s empirical program on human--AI
collaboration protocols, frictional AI, and technology dominance measurement.
\end{abstract}

\tableofcontents
\newpage

% ============================================================================
\section{Introduction}
\label{sec:introduction}
% ============================================================================

A convergence of independent research programs has established that heavy use of
large language models (LLMs) produces measurable cognitive degradation across
multiple domains. The Microsoft/CMU study ($n=319$) found that approximately 40\%
of AI-assisted tasks involved zero critical thinking
\citep{lee2025critical}. The MIT Media Lab ($n=54$) demonstrated that LLM users
showed the weakest brain connectivity across all EEG frequency bands during
writing tasks \citep{kosmyna2025brain}. \citet{barcaui2025retention} documented a
16\% reduction in knowledge retention at 45 days ($d=0.68$). And the Lancet
reported a 20\% relative decline in adenoma detection rates among physicians after
AI exposure \citep{lancet2025deskilling}.

The common failure mode---dumping fragments into an LLM, accepting completions,
outsourcing email and analysis---is a textbook instance of what cognitive
scientists call the \textbf{generation effect in reverse}
\citep{slamecka1978generation}: by letting the machine generate, the user skips
the encoding process that builds durable skill.

What follows is a concrete protocol to reverse this, grounded in the specific
cognitive mechanisms at work and designed for practitioners who understand what an
LLM is doing under the hood.

% ============================================================================
\section{Theoretical Foundations}
\label{sec:foundations}
% ============================================================================

Eight independent research programs converge on the same conclusion. Understanding
the mechanisms matters because the protocol's design exploits them.

\subsection{The Generation Effect}

\citet{slamecka1978generation} established that information actively produced is
retained significantly better than information passively received. When a user
lets an LLM write their email, they are in the ``read'' condition---they get
fluent output but encode almost nothing. \citet{kosmyna2025brain} confirmed this
neurologically: LLM users could not quote their own essays minutes after writing
them, and their brain connectivity progressively weakened over four months of
AI-assisted writing.

\subsection{Desirable Difficulties}

\citet{bjork1994desirable} showed that conditions slowing apparent
performance---spacing, interleaving, retrieval practice, generation---are
precisely what produce durable learning. LLMs eliminate every single one. The
Wharton/PNAS randomised controlled trial \citep{bastani2025learning}
demonstrated this causally: students with unrestricted GPT-4 scored \textbf{17\%
lower} on unassisted exams than students with no AI at all, while perceiving no
learning deficit. Students given a scaffolded GPT tutor---one that asked
questions instead of giving answers---showed no harm. Deskilling is a design
problem, not an inevitable consequence of AI use.

\subsection{Automation Complacency}

\citet{parasuraman2010complacency} showed that automation complacency operates
through attention reallocation: when users trust the automation, they redirect
cognitive resources elsewhere. Critically, their meta-analysis showed this effect
\textbf{cannot be overcome by training, instructions, or awareness alone}---it
requires structural intervention in the interaction itself. This is why
willpower-based approaches (``I'll just be more careful'') reliably fail.

\subsection{Bainbridge's Ironies of Automation}

\citet{bainbridge1983ironies} identified the deepest problem: the better the
automation gets, the rarer and more unusual the situations requiring independent
human judgment become---and the less practised the human is at handling them.
Users need \emph{more} skill as AI improves, not less, because they will only be
called upon for the hard cases.

\subsection{Cabitza's Deskilling Research Programme}

\citet{cabitza2017unintended} launched a decade of empirical work showing how AI
decision support erodes the very expertise it claims to augment. Three concepts
from this programme are foundational for the Forge Protocol:

\paragraph{Semiotic deskilling.} Introduced in \citet{cabitza2021cobra}, this
refers to the erosion not just of cognitive ability but of \emph{interpretive
capacity}---the ability to read situations, perceive subtle cues, and make sense
of ambiguous information. When users stop writing their own emails, they lose not
just writing fluency but the ability to read tone, anticipate interpretation, and
notice when something is off. The skill atrophies at the semiotic level, not just
the mechanical one.

\paragraph{Epistemic sclerosis.} \citet{natali2025frictional} define this as the
rigidification of knowledge processes when decision support systems are blindly
trusted. In a team context, this is what happens when everyone fragment-dumps
into the same LLM---the collective capacity for original thought calcifies. This
extends deskilling from an individual phenomenon to an organisational one.

\paragraph{The deskilling taxonomy.} \citet{natali2025taxonomy} provide the most
comprehensive treatment, identifying four distinct types:
\emph{cognitive/technical} (loss of reasoning and procedural skill),
\emph{semiotic} (loss of interpretive capacity), \emph{social} (erosion of
communication and interpersonal skills), and \emph{moral} (reduction of ethical
judgment capacity). All four apply directly to LLM-assisted knowledge work.

\subsection{Empirical Findings on Collaboration Protocols}

In \citet{cabitza2023rams}, three human--AI collaboration approaches were tested
across 12 radiologists and 44 ECG readers:

\begin{description}[style=nextline]
    \item[``Rams'' (AI-first)] The AI provides its recommendation first, the
    human reviews. This is the standard LLM chat pattern---and it produced the
    worst outcomes. Clinicians anchored to incorrect AI suggestions, their
    independent reasoning collapsed, and automation bias dominated.

    \item[``Hounds'' (human-first)] The human commits to their own assessment
    first, then sees the AI's recommendation. This preserved independent judgment
    while still capturing AI's value. It directly validates the protocol's
    ``generate before you retrieve'' principle.

    \item[``White Boxes'' (XAI-enhanced)] Adding explanations to AI output did
    \emph{not} reliably improve decision quality---and in some cases worsened it.
    This is the \textbf{white-box paradox}.
\end{description}

\citet{cabitza2025five} extended this across radiology, ECG, and endoscopy,
finding that a \textbf{displacement protocol}---where human and AI work
completely independently, then results are combined---consistently achieved
87--89\% accuracy, outperforming AI-first and even human-first protocols.

\subsection{The XAI Halo Effect}

In \citet{cabitza2024explanations} (Best Paper Award, xAI 2024), it was
demonstrated that when AI output comes with persuasive-looking explanations,
users accept it \emph{more uncritically}, not less. A confident, well-structured
LLM response feels like understanding. It is not. It is recognition masquerading
as comprehension. This finding directly threatens naive approaches to
anti-deskilling: Claude's articulate, well-reasoned responses are themselves a
deskilling vector.

\subsection{Convergence and Homogenisation}

\citet{doshi2024creativity} showed AI-assisted outputs are \textbf{5\% more
similar to each other} than human-only outputs. If a user's emails, ideas, and
writing are starting to sound the same---or sound like the LLM---they have
outsourced their voice. Cabitza's concept of epistemic sclerosis extends this to
the team level.

% ============================================================================
\section{Design Principles}
\label{sec:principles}
% ============================================================================

Eleven principles are ordered by importance and grounded in specific mechanisms.
Each directly informs the protocol. Principles 1--7 derive from general cognitive
science; principles 8--11 derive specifically from Cabitza's empirical programme.

\begin{enumerate}
    \item \textbf{Generate before you retrieve.} The generation effect means
    producing your own version first---even a bad one---creates encoding that
    passive consumption cannot. This is the single highest-leverage intervention.
    Cabitza's ``Hounds'' protocol validates this empirically
    \citep{cabitza2023rams}.

    \item \textbf{Commit before you consult.} Croskerry's cognitive forcing
    strategies \citep{croskerry2003forcing} require the clinician to form a
    differential diagnosis \emph{before} looking at test results, because the act
    of commitment activates analytical reasoning.

    \item \textbf{Friction is the mechanism, not the obstacle.} Bjork's research
    proves that the difficulty you want to remove is exactly what builds the
    skill. \citet{cabitza2024frictional} formalise this as ``Frictional AI'' and
    \citet{cabitza2019programmed} introduce ``programmed inefficiencies.''

    \item \textbf{Separate thinking from executing.} Not all LLM interactions
    carry deskilling risk. Formatting, translation, boilerplate---these are
    execution tasks. Strategy, arguments, judgment calls---these are thinking
    tasks where the AI must be a sparring partner, not a ghostwriter.

    \item \textbf{Scaffold, don't substitute.} The Wharton/PNAS study showed
    that scaffolded AI preserved learning while unscaffolded AI destroyed it.
    \citet{cabitza2022open} formalise this as the ``Open, Multiple, Adjunct''
    design principles.

    \item \textbf{The AI critiques; you create.} Reversing the default
    flow---where the user creates and the AI evaluates---preserves both the
    generation effect and the user's ownership of the output.

    \item \textbf{Explain back before you use.} Retrieval practice shows that
    re-articulating information in one's own words is one of the strongest
    consolidation mechanisms.

    \item \textbf{Beware the explanation trap (the white-box paradox).}
    AI explanations can increase automation bias rather than decrease it
    \citep{cabitza2024explanations}.

    \item \textbf{Design for judicial interaction, not oracular consultation.}
    \citet{cabitza2025judicial} propose shifting from AI that delivers a single
    authoritative answer to AI that presents contrasting perspectives.

    \item \textbf{Measure technology dominance.} The AFOOT framework
    \citep{cabitza2023dominion} provides metrics for measuring how much AI
    influences decisions.

    \item \textbf{Monitor for convergence.} \citet{doshi2024creativity} showed
    AI-assisted outputs converge. If an entire organisation routes thinking
    through the same LLM, collective knowledge rigidifies.
\end{enumerate}

% ============================================================================
\section{The Forge Protocol: Four Modes of LLM Interaction}
\label{sec:modes}
% ============================================================================

The core framework defines four distinct modes. Each has explicit rules about what
the LLM does, what the user does, and when to use it. The metaphor: the user is
the smith, the AI is the anvil. The anvil shapes nothing by itself---it provides
resistance against which the smith does the work.

The four modes map directly onto Cabitza's empirically tested collaboration
protocols: Forge mode implements his judicial interaction paradigm; Anvil mode
implements the Hounds protocol; Crucible mode implements frictional AI; and
Executor mode is the only context where the Rams (AI-first) protocol is
acceptable.

\subsection{Mode 1 --- Forge (Thinking Partner)}

\textbf{When to use:} Developing ideas, making decisions, working through complex
problems, forming positions on ambiguous questions. This is the default mode for
anything involving judgment.

\textbf{What the LLM does:} Asks Socratic questions. Steelmans and attacks the
user's positions. Surfaces assumptions. Identifies what the user has not
considered. Never provides the answer, always provides the pressure. This
implements Cabitza's judicial AI paradigm---presenting contrasting evidence rather
than delivering a verdict.

\textbf{What the user does:} States their position first. Answers the LLM's
questions. Does the synthesis. Writes the conclusion.

\textbf{The white-box paradox warning:} The LLM's responses in Forge mode will
be articulate and well-reasoned. This is the XAI halo effect at work---the
quality of the questioning can itself become a form of intellectual outsourcing if
the user starts deferring to the framing of the questions rather than engaging
with the substance.

\paragraph{Questioning protocol (cognitive forcing functions).}
\begin{enumerate}
    \item \emph{Definitional:} ``What specifically do you mean by X?''
    \item \emph{Evidential:} ``What's your strongest evidence for that?''
    \item \emph{Adversarial:} ``What would someone who disagrees say, and why?''
    \item \emph{Implication:} ``If that's true, what else must be true?''
    \item \emph{Gap:} ``What haven't you considered yet?'' (list the gaps but do
    not fill them)
\end{enumerate}

\paragraph{Judicial protocol.} When the user shares a position or draft:
\begin{itemize}
    \item Present the case \textsc{for} the user's position (steelman it)
    \item Present the case \textsc{against} (identify the 2--3 most serious
    weaknesses)
    \item Present contrasting evidence on both sides---do not converge to a
    single recommendation
    \item Then ask: ``Given both sides, how do you weigh this? What's your
    ruling?''
\end{itemize}

\paragraph{Frictional AI --- programmed inefficiencies.}
\begin{itemize}
    \item If the user gives a fragment and expects completion, refuse. Ask them to
    write three full sentences first.
    \item If the user accepts the LLM's framing too quickly, flag it: ``You
    agreed with that fast. What's your independent reasoning?''
    \item Periodically withhold assessment and ask the user to predict what the
    LLM will say before it says it.
\end{itemize}

\subsection{Mode 2 --- Anvil (Critic and Editor)}

\textbf{When to use:} After the user has written a full draft of anything. The
LLM evaluates their work. It never rewrites; it identifies problems and asks the
user to fix them.

This mode implements Cabitza's empirically validated Hounds protocol
\citep{cabitza2023rams}: the human produces their assessment first, the AI
responds second. The key insight is that the order of interaction is not a minor
UX choice---it fundamentally determines whether the human retains or surrenders
cognitive autonomy.

\textbf{What the LLM does:} Grades the draft against specific criteria.
Identifies weaknesses in logic, tone, clarity, persuasiveness. Points to specific
lines. Asks forcing questions about choices the user made. Returns the draft with
annotations, not a replacement.

\textbf{What the user does:} Writes the entire first draft themselves, even if
rough. Reads the critique. Revises manually.

\paragraph{Rating dimensions.}
\begin{enumerate}
    \item \emph{Clarity:} Can the reader understand this on first pass?
    \item \emph{Precision:} Are claims specific and evidence-based?
    \item \emph{Structure:} Does it flow logically? Is the most important point
    first?
    \item \emph{Tone:} Is it appropriate for the audience and purpose?
    \item \emph{Persuasiveness:} Would this achieve its goal?
    \item \emph{Concision:} What can be cut without losing meaning?
\end{enumerate}

\subsection{Mode 3 --- Crucible (Idea Development)}

\textbf{When to use:} Brainstorming, exploring a problem space, developing
strategy. The user generates ideas; the LLM stress-tests them.

This mode implements Cabitza's frictional AI concept and his programmed
inefficiencies design pattern \citep{cabitza2019programmed}. The core insight is
that decision support systems should include deliberate friction features that
prevent automatic reliance.

The risk in LLM-assisted brainstorming is \textbf{premature convergence}: the LLM
generates plausible-sounding ideas, the user anchors on them, and the possibility
space collapses \citep{doshi2024creativity}.

\paragraph{The Crucible discipline.}
\begin{enumerate}
    \item \textbf{Solo generation} (5--15 minutes, no LLM). Generate your own
    ideas first. Aim for quantity over quality. At least 5 ideas, including ones
    you think are bad.
    \item \textbf{Commitment and ranking.} Before opening the LLM, rank your top
    3 ideas and write one sentence explaining why each matters.
    \item \textbf{Enter Crucible mode.} Share ideas with the LLM. Ask it to
    steelman, attack, map gaps, and identify hidden assumptions.
    \item \textbf{You fill the gaps.} The LLM identified blind spots but did not
    fill them. The user does the work.
    \item \textbf{Divergence check.} Ask: ``How is my current thinking different
    from what a generic smart person would come up with?''
\end{enumerate}

\subsection{Mode 4 --- Executor (Automation)}

\textbf{When to use:} Tasks where cognitive skill is not at stake: formatting,
translation, data transformation, boilerplate generation, summarising documents
already read, code scaffolding for well-understood patterns.

This is the only mode where Cabitza's Rams (AI-first) protocol is acceptable. His
research shows the Rams protocol leads to anchoring and automation bias when
judgment is involved---but for mechanical tasks, those biases are irrelevant
because there is no judgment to corrupt.

\textbf{What the LLM does:} Whatever you tell it. No friction. Full automation.

\textbf{The critical discipline:} Honestly classifying the task. If the user
catches themselves using Executor mode for an email that requires judgment, an
analysis that requires their opinion, or a document that carries their name---they
must stop and switch modes.

\subsection{Mode Selection Decision Rule}

Before every LLM interaction, the user should ask: \textbf{``Am I about to think,
or about to delegate thinking?''}

\begin{itemize}
    \item Does this task require my judgment, voice, or expertise? $\rightarrow$
    Forge or Anvil
    \item Am I developing or evaluating an idea? $\rightarrow$ Forge or Crucible
    \item Have I already written a draft? $\rightarrow$ Anvil
    \item Is this mechanical transformation of known inputs? $\rightarrow$
    Executor
    \item Am I too tired to think but the task requires thinking? $\rightarrow$
    Do not use the LLM. Rest, then do it yourself. Fatigue + LLM = maximum
    deskilling.
\end{itemize}

% ============================================================================
\section{Replacing Fragment-Dumping: The 3-Before-1 Rule}
\label{sec:fragment}
% ============================================================================

The failure mode of dumping fragments into the LLM and letting it complete is the
highest-priority behaviour to replace. In Cabitza's terms, this is the purest
form of semiotic deskilling: the user loses not just writing skill but the
interpretive capacity to formulate their own thoughts into communicable language.

\textbf{Old pattern:} Fragment $\rightarrow$ LLM completes $\rightarrow$ user
accepts $\rightarrow$ done.

\textbf{New pattern (the ``3-before-1'' rule):}
\begin{enumerate}
    \item \textbf{Write three full sentences yourself} before opening the chat.
    Not bullets. Not fragments. Ugly, rough, grammatically imperfect sentences
    that express your actual thought. This activates the generation effect.

    \item \textbf{State your intent in one sentence} at the top of the chat: ``I'm
    trying to [specific goal] for [specific audience] because [specific reason].''
    This is Croskerry's cognitive forcing---committing to a frame before receiving
    input.

    \item \textbf{Paste your rough sentences} and ask the LLM (in Forge mode):
    ``What's unclear? What am I assuming? What's the strongest objection?'' Do
    \emph{not} ask it to ``expand'' or ``flesh out'' or ``turn this into.''

    \item \textbf{You write the full version yourself,} incorporating the LLM's
    challenges. Then switch to Anvil mode for critique.
\end{enumerate}

If the user cannot write three sentences about the topic, that is diagnostic
information. It means they do not yet know what they think---and that is exactly
when they most need to do the thinking themselves, not outsource it.

% ============================================================================
\section{Email Protocol: Reclaiming the Draft}
\label{sec:email}
% ============================================================================

Email is the highest-frequency deskilling vector because it is low-stakes enough
to feel like Executor territory but high-frequency enough to erode writing
instincts through sheer volume. Cabitza's deskilling taxonomy helps classify
what is at risk: even a ``simple'' email involves \emph{social} skill (reading the
relationship), \emph{semiotic} skill (choosing register and tone), and sometimes
\emph{moral} skill (deciding how candid to be). Only truly templated messages are
purely technical.

\textbf{For routine/low-stakes emails} (scheduling, acknowledgements, simple
requests): Executor mode is fine.

\textbf{For any email that requires judgment} (persuasion, bad news,
cross-functional requests, anything to leadership):

\begin{enumerate}
    \item \textbf{Capture intent} (30 seconds, no LLM): To, Goal, Subtext,
    Constraint.
    \item \textbf{Write the ugly draft} (2--5 minutes, no LLM). Allow it to be
    clumsy. The point is generation, not quality.
    \item \textbf{Critique via Anvil.} The LLM rates and challenges your draft.
    It does not rewrite.
    \item \textbf{You revise.}
    \item \textbf{Optional final pass:} ``Read this as [recipient]. What's your
    emotional reaction? What might you misread?''
\end{enumerate}

% ============================================================================
\section{The User's Constitution: Twelve Rules for Cognitive Sovereignty}
\label{sec:constitution}
% ============================================================================

These personal disciplines create structural barriers that do not depend on
willpower in the moment---which, per \citet{parasuraman2010complacency}, is the
only intervention that actually works against automation complacency.

\begin{enumerate}
    \item \textbf{Always write before you prompt.} For any thinking task, produce
    at least three rough sentences before engaging the LLM. No exceptions.

    \item \textbf{Never accept the first completion.} If the LLM generates text
    you plan to use, you must revise it substantively---change the structure, not
    just word-swap.

    \item \textbf{Explain back before you use.} After the LLM gives a complex
    response, close the chat and articulate the key insight in your own words
    before incorporating it.

    \item \textbf{Name the mode before you type.} Every LLM session starts with a
    conscious decision: Forge, Anvil, Crucible, or Executor.

    \item \textbf{Maintain a ``done without AI'' list.} Keep a running log of
    things completed without AI assistance each week.

    \item \textbf{One AI-free writing block per week.} Minimum 45 minutes of
    writing with no AI access.

    \item \textbf{Read your sent emails.} Once a week, reread 10 sent emails. Do
    they sound like you, or like the LLM?

    \item \textbf{Don't use AI when tired.} Fatigue is when automation
    complacency is strongest.

    \item \textbf{Preserve struggle on things that matter to your identity.}
    Protect skills across all four dimensions of
    Cabitza's taxonomy---technical, semiotic, social, and moral.

    \item \textbf{Track your prompts.} Count the ratio of ``write this for me''
    to ``challenge me on this.'' The ratio is your dependency score
    \citep{cabitza2023dominion}.

    \item \textbf{When the LLM is clearly better, learn why.} Identify the
    specific technique. The goal is skill transfer, not skill replacement.

    \item \textbf{Treat discomfort as signal, not friction to eliminate.} The
    moment of not-knowing-what-to-say-next is exactly when the user most needs to
    stay with it.
\end{enumerate}

% ============================================================================
\section{Self-Audit Practices}
\label{sec:audit}
% ============================================================================

Deskilling is insidious precisely because it feels like efficiency.
\citet{cabitza2023dominion} provide the theoretical basis for measurement; the
following are the practical detection methods.

\paragraph{Weekly canary check (5 minutes).}
\begin{itemize}
    \item Did I write anything substantial ($>$200 words) without AI this week?
    \item Can I recall the main argument of the last document I produced?
    \item Did I change my mind about anything this week, or did I just have the AI
    validate my initial position?
    \item Did any LLM output surprise me?
    \item How many times did I override or substantially revise AI output?
\end{itemize}

\paragraph{Monthly stress test (30 minutes).} Write a one-page analysis, memo, or
argument on a familiar topic with no AI assistance. Time yourself. Compare
quality and speed to six months ago. This tests cognitive/technical deskilling.

\paragraph{Monthly semiotic test (15 minutes).} Read 5 emails received this week.
For each, write down: What is this person actually trying to say? What is the
subtext? This tests semiotic skill---the interpretive capacity that atrophies when
one stops engaging closely with language.

\paragraph{Quarterly dependency audit.} Review LLM conversation history for the
past month. Categorise each session as Forge/Anvil/Crucible/Executor. If Executor
is $>$70\% and you are doing knowledge work, you are likely deskilling. The
healthy ratio for a senior knowledge worker is roughly \textbf{40\%
Forge/Crucible, 30\% Anvil, 30\% Executor}.

\paragraph{Red flags that demand immediate correction:}
\begin{itemize}
    \item Opening a chat window before forming any thought about the task
    \item Inability to write a professional email without AI assistance
    \item Noticeably less structure or clarity in unassisted writing
    \item Inability to explain your own AI-assisted work in meetings
    \item Accepting AI-generated reasoning without independently agreeing
    \item Colleagues' writing all sounding the same (organisational epistemic
    sclerosis)
\end{itemize}

% ============================================================================
\section{Where Cognitive Forcing Is Not Worth It}
\label{sec:triage}
% ============================================================================

Applying this protocol to everything is a recipe for burnout and abandonment. The
point is not to maximise difficulty---it is to preserve difficulty where it
matters.

\textbf{Full cognitive forcing} (Forge/Anvil/Crucible): Strategy documents,
important emails, performance reviews, analysis that carries your name, decision
memos, idea development, anything you will need to defend in a meeting.

\textbf{Light cognitive forcing} (quick self-draft $\rightarrow$ AI critique):
Medium-stakes emails, slide decks, documentation, internal communications where
tone matters.

\textbf{No cognitive forcing} (Executor, full automation): Calendar coordination,
templated messages, data formatting, translation, code boilerplate, summarising
documents already read, research summarisation for topics where you will do your
own synthesis.

\paragraph{The honest tradeoff.} This protocol will make you slower on tasks
where you are currently fast because the AI is doing the work. That is the cost.
The benefit is that in six months, you can still think without the AI---and you
will be better at catching its mistakes because you have maintained the cognitive
capacity to evaluate its output independently.

The deepest finding from the Wharton study \citep{bastani2025learning} is that
deskilling and productivity gains coexist invisibly. The students who were getting
worse scored better on every assisted task. They never noticed. Neither will
you---unless you build the measurement infrastructure to catch it.

% ============================================================================
\section*{Acknowledgements}
% ============================================================================

This protocol draws extensively on the empirical programme of Federico Cabitza and
collaborators at the University of Milano-Bicocca. Their decade of work on
human--AI collaboration, frictional AI design, and technology dominance
measurement provides the scientific foundation for the interaction patterns
described here.

% ============================================================================
\bibliographystyle{plainnat}
\bibliography{references}

\end{document}
